\section*{第 \ref{ch:b3:microcanonical-ensemble} 章——微正则系综}

\begin{solution}{exo:b3:microcanonical-ensemble:1}
(a) $1, 4, 6, 4, 1$。(b) 以 $16$ 为分母：$6.25\%$、$25\%$、$37.5\%$、
$25\%$、$6.25\%$。(c) 依次为 $S/k_{\text{B}} = 0$、$\ln4$、$\ln6$、$\ln4$、
$0$。(d) $n = 2$，仅占 $37.5\%$ 的概率：在
$N = 4$ 时，“平衡”仅为相对多数——使热力学成立的清晰性
须以大 $N$ 换取。
\end{solution}

\begin{solution}{exo:b3:microcanonical-ensemble:2}
(a) 当 $N = 10$ 时为 $13.0$，与 $15.1$ 相比（偏差 $14\%$）；当 $N = 100$ 时：
$360.5$ 与 $363.7$ 相比（偏差 $0.9\%$）。(b) 略去的项为
$\tfrac12\ln(2\pi N)$，相对误差 $\sim\ln N/N$。(c) 对三个阶乘应用
斯特林公式。(d) 当 $N = 10^{23}$ 时，略去项
为 $\sim26$，与 $N\ln2 \sim 10^{23}$ 相比，
比主项低了二十二个数量级。
\end{solution}

\begin{solution}{exo:b3:microcanonical-ensemble:3}
(a) $\ln(\Omega_1\Omega_2) = \ln\Omega_1 + \ln\Omega_2$。(b)
$\Delta S = k_{\text{B}}\ln2$——一个粒子的
可用体积加倍。(c) $S = 156\,k_{\text{B}} = \qty{2.2e-21}{J/K}$。(d)
热力学熵带有 $10^{23}$ 量级的因子：地球上所有
赌场的洗牌，相比一口
温热空气，几乎不可察觉。
\end{solution}

\begin{solution}{exo:b3:microcanonical-ensemble:4}
(a) $\Delta S = N_{\text{A}}k_{\text{B}}\ln2 = R\ln2 =
\qty{5.76}{J/K}$。(b) $2^{-N_{\text{A}}} \sim 10^{-1.8\times
10^{23}}$。(c) 即便每隔 \qty{e-10}{s} 抽样一次构型，
预期等待时间远超 \qty{4e17}{s}，其倍数的
\emph{指数}有二十三位数字。(d) 并非不可能——
无从见证：定律原则上是概率性的，
而在任何包含有限耐心观测者的世界中则是绝对的。
\end{solution}

\begin{solution}{exo:b3:microcanonical-ensemble:5}
(a) 对乘积取对数。(b) 将 $V, N$ 加倍，
保持 $E/N$ 不变：若无 $N!$，$S$ 将多出一个 $N\ln2$ 型的超额量
——吉布斯佯谬。(c) 对 $\ln N!$ 用斯特林公式，
可得 $V/N$ 与 $E/N$ 的组合。(d) 普朗克相格 $h^3$ 确定了
熵的零点；粒子的全同性确定了 $N!$——两个
隐藏于经典热力学中的量子输入。
\end{solution}

\begin{solution}{exo:b3:microcanonical-ensemble:6}
(a) $1/T = \tfrac32 Nk_{\text{B}}/E$。(b) $P/T =
Nk_{\text{B}}/V$：由计数得到的理想气体定律。(c)
$S$ 为常数：$\ln(V/N) + \tfrac32\ln(E/N)$ 固定，于是 $VE^{3/2}$ 为常数；
由 $E \propto T$ 得 $TV^{2/3}$ 为常数，
与 $PV^{5/3}$ 等价。(d) 气体定律、能量与温度的关系以及
绝热指数——本科一年级分册里的三大经验支柱——
如今都已成为定理。
\end{solution}

\begin{solution}{exo:b3:microcanonical-ensemble:7}
(a) 对称性（各块相等）。(b) 每个 $\ln\Omega \propto
\tfrac32 N\ln E_i$：在均分点附近展开，
得到关于 $x$ 的高斯分布，方差为 $\sim1/N$。(c) $x_{\text{rms}} \sim
\num{e-11}$。(d) $\Delta T/T \sim x_{\text{rms}}$：
纳纳开尔文——涨落确实存在，并且在
实验室尺度上小到无法测量。
\end{solution}

\begin{solution}{exo:b3:microcanonical-ensemble:8}
(a) $\Omega = \binom Nn$，$S = k_{\text{B}}\ln\binom Nn$。(b)
$1/T = \partial S/\partial E = (k_{\text{B}}/\epsilon)[\ln(N-n) -
\ln n]$；解出 $n/N$。(c) $E(T)$ 从 $0$ 升至饱和值，
即 $N\epsilon/2$；$C = \dd E/\dd T$ 在低 $T$ 时趋于零
（此时无量子可被吸收），高 $T$ 时亦趋于零（两能级同等占满，
已无可吸收）：二者之间出现一个峰——这是任何双能级粒子数分布的量热
指纹，可用于探测
固体中的磁性杂质。(d) 它是热致双能级占据
——并且，令 $\epsilon \to E - \mu$，它将原样重现为
费米–狄拉克分布。
\end{solution}

\begin{solution}{exo:b3:microcanonical-ensemble:9}
(a) $1/T \propto \ln[(N-n)/n] < 0$，当且仅当 $n > N/2$ 时成立。(b)
$T = \epsilon/[k_{\text{B}}\ln(2/3)] = -\qty{2.9}{K}$。(c) 一个
负 $T$ 的物体通过失去能量而\emph{增加}熵，
正 $T$ 的物体通过获得能量而增加熵：两者的计数都随能量
由负 $\to$ 正流动时均增长——倒置系统是
通用施主，即比一切都热。(d) 在能量无界的情形下，
$\Omega(E)$ 永不回落：$\partial S/\partial E$
在任何有限能量下仍保持为正。
\end{solution}

\begin{solution}{exo:b3:microcanonical-ensemble:10}
(a) 每种气体的可及体积均加倍：$\Delta S =
2Nk_{\text{B}}\ln2$。(b) 合并后各群体的 $N!$
恰好吸收了表观增益：$\Delta S = 0$。(c) 熵的
增加正是\emph{有序性的丧失}：没有任何宏观量
发生变化，但重新分离现在需要做功——这部分增加
被储存为一笔未来的账单。(d) $W_{\min} = T\Delta S =
2Nk_{\text{B}}T\ln2$；同位素分离工厂（铀
离心级联、氦-3 回收）的实际耗功
往往是这一热力学下限的许多倍。
\end{solution}

\begin{solution}{exo:b3:microcanonical-ensemble:11}
(a) $\lambda_T = \qty{1.6e-11}{m}$。(b) $V/N = k_{\text{B}}T/P =
\qty{4.1e-26}{m^3}$：$V/N\lambda_T^3 = \num{1.0e7}$，
于是 $S/Nk_{\text{B}} = \ln(\num{1.0e7}) + 2.5 = 18.6$。(c) 与
量热测得的 $18.6$ 精确吻合。(d) 绝对熵包含 $h$
（通过 $\lambda_T$）与 $N!$ 项：来自蒸汽时代的热测量
证实到三位有效数字的作用量子，以及
原子的全同性。
\end{solution}

\begin{solution}{exo:b3:microcanonical-ensemble:12}
(a) $v = (\qty{550}{nm})^3 = \qty{1.7e-19}{m^3}$：$\bar n =
\num{4.2e6}$，$\delta_{\text{方均根}} = 1/\sqrt{\bar n} \approx
\num{5e-4}$。(b) 完全均匀的介质
不会向侧向散射任何光（见\cref{pb:b3:perturbation-theory:1}，第三部分）：这些
在每个光学元胞中不可约的 $\sqrt N$ 涨落
正是使天空得以存在的无序。(c) 取 $\bar n = 100$：$v
\approx (\qty{16}{nm})^3$。(d) 天空呈蓝色，因为空气由可数的分子构成，
其数目以 $\sqrt N$ 的量级涨落——
爱因斯坦正是借此论证分子确实存在。
\end{solution}

\begin{solution}{pb:b3:microcanonical-ensemble:1}
\textbf{1.} $n_\pm = \tfrac12(N \pm x/b)$；$\Omega =
N!/n_+!\,n_-!$。
\textbf{2.} 当 $x = 0$ 时：自由链卷曲；完全伸展时
微观态数 $\Omega = 1$。
\textbf{3.} 对二项式取斯特林近似，展开到关于
$x/Nb$ 的二阶：即所述的高斯熵。
\textbf{4.} 由于 $E$ 与 $x$ 无关，任何使
$x$ 变小的拉力都只能来自计数：按构造即为熵效应。
\textbf{5.} 线团尺度 $\sim b\sqrt N = \qty{50}{nm}$，对比 $Nb =
\qty{5}{\micro m}$：自由链比其轮廓长度
短一百倍——几乎完全卷曲。
\textbf{6.} $N$，即交联点之间的链段数：
硫化缩短了有效链，使
网络变硬。
\textbf{7.} $f = -T\,\partial S/\partial x =
k_{\text{B}}Tx/Nb^2$。
\textbf{8.} $k = k_{\text{B}}T/Nb^2 = \num{4.14e-21}/(\num{e4}
\times \num{2.5e-19}) = \qty{1.7e-6}{N/m}$，每条链。
\textbf{9.} $k \propto T$：在固定载荷下加热时，橡皮筋
变硬并\emph{收缩}，抬起载荷；钢弹簧
则只是变软，略有下垂。
\textbf{10.} $f(Nb/2) = k_{\text{B}}T/2b \approx \qty{4}{pN}$；
线性定律在 $f$ 趋近 $k_{\text{B}}T/b \approx
\qty{8}{pN}$ 时失效，此时链接近完全伸展。
\textbf{11.} 并联的链约 $\sim\num{e14}$ 条，而
沿一厘米方向串联约 $\sim\num{e5}$ 个线团长度：
$k_{\text{橡皮筋}} \sim k_{\text{链}} \times \num{e14}/\num{e5}
\sim \num{e2}$--$\num{e3}\,\unit{N/m}$——这正是
橡皮筋给人的熟悉手感。
\textbf{12.} 橡胶的劲度系数\emph{正是} $k_{\text{B}}T$ 每个
构型：热运动越剧烈，回弹越强；金属的劲度系数
则为键能，非简谐热运动会使其松弛。
\textbf{13.} 拉伸减少构型数；若快速进行
（与外界无熵交换），失去的构型
熵必须以热熵的形式重现：橡皮筋升温——
\qty{1}{K} 量级，嘴唇可以察觉。
\textbf{14.} 放松使构型恢复；热量
核算退还了这部分差额：橡皮筋降温。
\textbf{15.} 重物\emph{上升}：加热时的收缩
即熵效应的特征（能量弹性材料则膨胀）。
\textbf{16.} 钢弹簧略有伸长（热
膨胀），其模量下降：符号相反——属于能量弹性，
而非计数弹性。
\textbf{17.} 受热的辐条收缩，把轮缘的质量
偏离中心、拉向较冷一侧；重力对失衡的
车轮施加力矩；每根辐条在转离时重新松弛：一台热机
其工作物质就是熵本身。
\textbf{18.} $k_{\text{B}}T/b = \num{4.14e-21}/\num{e-7} =
\qty{4e-14}{N} = \qty{0.04}{pN}$：皮牛及以下——
恰在光镊的范围内，因此
单个 DNA 分子的力–伸长定律可以逐点追踪（并与
第 23 条中的修正相吻合）。
\textbf{19.} 纯粹来自计数：在伸长 $x$ 处
的微观态比 $x - \dd x$ 处的少，热运动
使链向微观态更多的方向漂移；所谓“力”
即 $T$ 乘以熵梯度。
\textbf{20.} 当 $T = 0$ 时，没有任何构型被探索：
熵力随驱动它的热运动一同消失——
弹性将纯粹是能量性的，橡胶
像一根软塌的线。
\textbf{21.} 气体：$P = T\,\partial S/\partial V$，其中 $S \ni
Nk_{\text{B}}\ln V$；链：$f = -T\,\partial S/\partial x$，
采用高斯计数——两者都是计数在某一
约束下的推动，并以 $T$ 计价。
\textbf{22.} 吹风机的热量进入橡皮筋；部分
在等温收缩中转化为功，账目由
进出的熵来平衡——一台微型热机，热力学第一定律
依然成立。
\textbf{23.} 有限伸长性（真实的力在接近
完全拉伸时发散——蠕虫状链修正）以及
应变诱导的结晶（取向使链有序，
产生滞后和超出理想模型的放热）。
\textbf{24.} $k_{\text{B}}T/\unit{nm} = \qty{4.1}{pN}$：这是
热能与纳米位移恰好平衡的力——
分子马达、折叠蛋白与被拉伸的 DNA
都在其一个数量级之内运作。
\textbf{25.} 一万条盲链，计数可得：胡克定律，
$k = k_{\text{B}}T/Nb^2$，拉伸时升温，重物被
热空气举起，单个分子在 \qty{0.04}{pN} 处服从——
熵作为一种力在起作用。
\end{solution}
